\newpage 

\section*{What changes when $G=1$?}

\begin{itemize}
    \item Each group contains just one response $Y^{(b)} \implies$  leave-one-out advantage is no longer defined, so we replace it with
    \begin{equation*}
         A(x, y)=r_\phi(x, y)-V(x),   
    \end{equation*}
    \item Importance ratio becomes 
    \begin{equation}
    \rho_n^{(b)}=\frac{\pi_{\theta_n}\left(Y^{(b)} \mid x^{(b)}\right)}{\pi_\theta\left(Y^{(b)} \mid x^{(b)}\right)} 
    \end{equation}
    and $\bar{\rho}_n^{(b)}=\rho_n^{(b)} .$
    \item The GMC ratio simplifies to
    \begin{equation*}
        \hat{\rho}_n^{(b)}=\min \left\{\rho_n^{(b)}, 1+\epsilon\right\} .
    \end{equation*}
    \item Group-level gradient estimator
    \begin{equation*}
        U_{n, b}=\frac{1}{T} \min \left\{\rho_n^{(b)}, 1+\epsilon\right\} \nabla \log \pi_{\theta_n}\left(Y^{(b)} \mid x^{(b)}\right) A\left(x^{(b)}, Y^{(b)}\right).
    \end{equation*}
\end{itemize}

\subsection*{Biases}